% Canonical derivation source for the weak-anisotropy unfolding of the ZGV ring.
% This file is a LaTeX fragment and is included by the Supplementary Information.

\section{Weak-anisotropy normal form and Morse splitting}
\label{sec:anisotropic-morse-derivation}

This section closes the coefficients of the weak-anisotropy normal form
directly from the Hermitian elastic eigenproblem.  We first derive the
eigenvalue sensitivity under arbitrary normalization.  We then specialize
it to the fixed-density plate and differentiate the mode in a fixed mass
gauge.  A separate tensor calculation shows why the registered cubic family
has only a constant and a fourfold term at first order.  Finally, a radial
implicit-function reduction and the exact Cartesian Hessian give the number
and Morse type of the split critical points.

Throughout, a dagger denotes the Hermitian adjoint and
\(\lambda=\omega^2\).  The in-plane polar basis is
\begin{equation}
  \mathbf n(\theta)=(\cos\theta,\sin\theta,0),
  \qquad
  \mathbf t(\theta)=(-\sin\theta,\cos\theta,0),
  \qquad
  \mathbf z=(0,0,1).
  \label{eq:anisotropic-polar-basis}
\end{equation}
The radial coordinate is denoted by \(k>0\), and
\(q=k-k_0\) measures distance from the isotropic ZGV circle.

\subsection{General Hermitian-pencil sensitivity}

Let a real parameter \(p\) enter a Hermitian generalized eigenproblem
\begin{equation}
  K(p)\mathbf u(p)=\lambda(p)M(p)\mathbf u(p),
  \qquad
  K(p)^\dagger=K(p),
  \qquad
  M(p)^\dagger=M(p)>0.
  \label{eq:general-hermitian-pencil}
\end{equation}
Assume that \(\lambda\) is simple and let
\begin{equation}
  N:=\mathbf u^\dagger M\mathbf u>0.
  \label{eq:arbitrary-mass-norm}
\end{equation}
No normalization has yet been imposed.  Differentiating
Eq.~\eqref{eq:general-hermitian-pencil} gives
\begin{equation}
  (K_p-\lambda_pM-\lambda M_p)\mathbf u
  +(K-\lambda M)\mathbf u_p=\mathbf0.
  \label{eq:differentiated-general-pencil}
\end{equation}
Left multiplication by \(\mathbf u^\dagger\) removes the last term because
\(\mathbf u^\dagger(K-\lambda M)=\mathbf0^\dagger\).  Therefore
\begin{equation}
  \lambda_p
  =\frac{\mathbf u^\dagger(K_p-\lambda M_p)\mathbf u}{N},
  \qquad
  \omega_p
  =\frac{\mathbf u^\dagger(K_p-\lambda M_p)\mathbf u}
         {2\omega N}.
  \label{eq:general-frequency-sensitivity}
\end{equation}
Both the numerator and denominator acquire the same factor \(|c|^2\) under
\(\mathbf u\mapsto c\mathbf u\), for any nonzero complex \(c\).  Thus
Eq.~\eqref{eq:general-frequency-sensitivity} is valid with arbitrary
normalization and is also phase invariant.  Hermiticity makes its numerator
real.

The present perturbation changes stiffness but not density.  Moreover, the
plate mass form is independent of the in-plane wave vector.  Hence the
fixed-density specialization is
\begin{align}
  M_k=M_\varepsilon=0,
  \qquad&
  \lambda_k=\frac{\mathbf u^\dagger K_k\mathbf u}{N},
  \notag\\
  \lambda_\varepsilon
  =\frac{\mathbf u^\dagger K_\varepsilon\mathbf u}{N},
  \qquad&
  V(k,\theta):=\omega_\varepsilon
  =\frac{\lambda_\varepsilon}{2\omega}.
  \label{eq:fixed-mass-first-sensitivities}
\end{align}
Here every derivative is taken on the tracked simple branch, and
\(K_\varepsilon\) is the stiffness derivative per unit \(\varepsilon\).

For completeness, the radial derivative of \(V\) can also be written before
choosing a normalization.  At fixed mass define
\begin{align}
  S_\varepsilon:=\mathbf u^\dagger K_\varepsilon\mathbf u,
  \qquad&
  (S_\varepsilon)_k
  =\mathbf u_k^\dagger K_\varepsilon\mathbf u
   +\mathbf u^\dagger K_{k\varepsilon}\mathbf u
   +\mathbf u^\dagger K_\varepsilon\mathbf u_k,
  \notag\\
  N_k=\mathbf u_k^\dagger M\mathbf u
      +\mathbf u^\dagger M\mathbf u_k.&
  \label{eq:arbitrary-normalization-radial-ingredients}
\end{align}
Applying the quotient rule to
\(V=S_\varepsilon/(2\omega N)\), and using
\(\omega_k=\lambda_k/(2\omega)\), yields
\begin{equation}
  \partial_kV
  =\frac{(S_\varepsilon)_k-S_\varepsilon N_k/N}{2\omega N}
   -\frac{S_\varepsilon\lambda_k}{4\omega^3N}.
  \label{eq:arbitrary-normalization-B}
\end{equation}
The combination in the first numerator cancels the derivative of any
\(k\)-dependent scale or phase of \(\mathbf u\).  Equation
\eqref{eq:arbitrary-normalization-B} is therefore normalization invariant,
as required by the left-hand side.

\subsection{Mass gauge, augmented solve, and the complete radial coefficient}

For computation and for a compact reduced-resolvent formula, choose the
mass-normalized eigenvector and its parallel-transport phase,
\begin{equation}
  \mathbf u^\dagger M\mathbf u=1,
  \qquad
  \mathbf u^\dagger M\mathbf u_k=0.
  \label{eq:mass-normalized-parallel-gauge}
\end{equation}
The real part of the second condition follows by differentiating the first;
the phase of \(\mathbf u\) is chosen to set its imaginary part to zero.  Put
\(A:=K-\lambda M\).  Equation
\eqref{eq:differentiated-general-pencil}, specialized to \(p=k\), together
with the gauge is equivalent to the augmented system
\begin{equation}
  \begin{pmatrix}
    K-\lambda M & M\mathbf u\\
    \mathbf u^\dagger M & 0
  \end{pmatrix}
  \begin{pmatrix}\mathbf u_k\\ \alpha\end{pmatrix}
  =
  \begin{pmatrix}
    -(K_k-\lambda_kM)\mathbf u\\ 0
  \end{pmatrix}.
  \label{eq:augmented-mode-derivative}
\end{equation}
For an exact simple eigenpair the multiplier is \(\alpha=0\).  The augmented
matrix is nonsingular: a homogeneous solution first gives
\(\alpha=0\) after left multiplication of its upper row by
\(\mathbf u^\dagger\); simplicity then gives
\(\mathbf x=c\mathbf u\), and the lower row forces \(c=0\).

Let \(\{(\lambda_n,\mathbf u_n)\}\) be a complete set of generalized
eigenpairs with
\(\mathbf u_m^\dagger M\mathbf u_n=\delta_{mn}\), and let the selected
mode have index zero.  Taking the inner product of the differentiated
equation with \(\mathbf u_n\), for \(n\ne0\), gives
\begin{align}
  (\lambda_n-\lambda)
  \mathbf u_n^\dagger M\mathbf u_k
  &=-\mathbf u_n^\dagger
       (K_k-\lambda_kM)\mathbf u                                      \notag\\
  &=-\mathbf u_n^\dagger K_k\mathbf u.
  \label{eq:mode-derivative-coefficient}
\end{align}
The gauge removes the component along \(\mathbf u\), so
\begin{equation}
  \mathbf u_k
  =\sum_{n\ne0}
    \mathbf u_n
    \frac{\mathbf u_n^\dagger K_k\mathbf u}{\lambda-\lambda_n}.
  \label{eq:mode-derivative-modal-sum}
\end{equation}
Equivalently, define the reduced resolvent
\begin{equation}
  \mathcal R_\lambda
  :=\sum_{n\ne0}
    \frac{\mathbf u_n\mathbf u_n^\dagger}{\lambda_n-\lambda}.
  \label{eq:reduced-resolvent}
\end{equation}
Then
\(
  \mathbf u_k=-\mathcal R_\lambda
  (K_k-\lambda_kM)\mathbf u
\).
The modal sum and the augmented system are two representations of the same
gauge-fixed derivative.  On a compact parameter set, a uniform positive
eigengap together with uniform coercivity and boundedness of the mass matrix
bounds \(\mathcal R_\lambda\) in the associated operator norms.  Here the
mass matrix is fixed and positive definite.  Loss of either hypothesis is
addressed below as a theorem boundary.

In the gauge \eqref{eq:mass-normalized-parallel-gauge}, differentiation of
\(\lambda_\varepsilon=\mathbf u^\dagger K_\varepsilon\mathbf u\)
gives the full mixed eigenvalue derivative
\begin{equation}
  \lambda_{k\varepsilon}
  =\mathbf u^\dagger K_{k\varepsilon}\mathbf u
   +2\operatorname{Re}
     (\mathbf u^\dagger K_\varepsilon\mathbf u_k).
  \label{eq:complete-mixed-eigenvalue-derivative}
\end{equation}
Consequently the radial coefficient in the frequency normal form is
\begin{align}
  B(\theta)
  :=\left.\partial_k\partial_\varepsilon\omega
    \right|_{(k_0,\theta,0)}
  &=\left.
    \left(
      \frac{\lambda_{k\varepsilon}}{2\omega}
      -\frac{\lambda_\varepsilon\lambda_k}{4\omega^3}
    \right)\right|_{(k_0,\theta,0)}                                  \notag\\
  &=\left.
    \left[
      \frac{\mathbf u^\dagger K_{k\varepsilon}\mathbf u}{2\omega}
      +\frac{\operatorname{Re}
        (\mathbf u^\dagger K_\varepsilon\mathbf u_k)}{\omega}
      -\frac{\lambda_\varepsilon\lambda_k}{4\omega^3}
    \right]\right|_{(k_0,\theta,0)}.
  \label{eq:complete-radial-frequency-sensitivity}
\end{align}
The three displayed contributions are, respectively, the explicit mixed
stiffness derivative, the differentiated-mode contribution, and the
frequency-denominator contribution.  At the exact isotropic ZGV point,
\(\lambda_k=2\omega\omega_k=0\), so the last contribution evaluates to
zero there.  It remains part of the general identity and is required away
from the exact ZGV radius.  In particular, \(B\) is not a fixed-mode strain
integral.

\subsection{Rotation and complete contraction of the cubic perturbation}

Let \(\{\mathbf a_1,\mathbf a_2,\mathbf a_3\}\) be the orthonormal cubic
axes.  A cubic elasticity tensor can be decomposed exactly as
\begin{align}
  C_{ijkl}
  ={}&C_{12}\delta_{ij}\delta_{kl}
   +C_{44}(\delta_{ik}\delta_{jl}+\delta_{il}\delta_{jk})
  \notag\\
  &+\Delta_C\sum_{\alpha=1}^{3}
      a_{\alpha i}a_{\alpha j}a_{\alpha k}a_{\alpha l},
  \qquad
  \Delta_C=C_{11}-C_{12}-2C_{44}.
  \label{eq:cubic-tensor-decomposition}
\end{align}
Indeed, the \(iiii\), \(iijj\), and \(ijij\) components of the right-hand
side are \(C_{12}+2C_{44}+\Delta_C=C_{11}\), \(C_{12}\), and
\(C_{44}\), respectively, and cubic symmetry generates all remaining
nonzero components.

The registered bulk-preserving family is
\begin{equation}
  C_{11}=\lambda+2\mu+\frac{2}{3}\varepsilon\delta,
  \qquad
  C_{12}=\lambda-\frac{1}{3}\varepsilon\delta,
  \qquad
  C_{44}=\mu.
  \label{eq:registered-cubic-family}
\end{equation}
It has
\begin{equation}
  C_{12}^{(1)}=-\frac{\delta}{3},
  \qquad
  C_{44}^{(1)}=0,
  \qquad
  \Delta_C^{(1)}
  :=\left.\partial_\varepsilon\Delta_C\right|_0=\delta,
  \qquad
  \Delta_C(\varepsilon)=\varepsilon\delta.
  \label{eq:cubic-family-derivatives}
\end{equation}
Thus Eq.~\eqref{eq:cubic-tensor-decomposition} gives
\begin{equation}
  C^{(1)}_{ijkl}
  =-\frac{\delta}{3}\delta_{ij}\delta_{kl}
   +\delta\sum_{\alpha=1}^{3}
      a_{\alpha i}a_{\alpha j}a_{\alpha k}a_{\alpha l}.
  \label{eq:registered-cubic-perturbation-tensor}
\end{equation}
The cubic stability inequalities reduce to
\(\mu>0\), \(2\mu+\varepsilon\delta>0\), and
\(3\lambda+2\mu>0\).  Thus, for the registered \(\delta>0\), stability
requires \(\varepsilon>-2\mu/\delta\); the weak signed interval used below is
chosen as a compact neighbourhood of zero inside this bound.  No assertion is
made for arbitrarily large negative \(\varepsilon\).

For a [001]-cut plate, take the crystallographic axes to be
\(\mathbf a_1=\mathbf e_x\),
\(\mathbf a_2=\mathbf e_y\), and
\(\mathbf a_3=\mathbf e_z\).  The rotation whose columns are
\((\mathbf n,\mathbf t,\mathbf z)\) is
\begin{equation}
  R(\theta)=
  \begin{pmatrix}
    \cos\theta&-\sin\theta&0\\
    \sin\theta& \cos\theta&0\\
    0&0&1
  \end{pmatrix},
  \qquad
  \widetilde C^{(1)}_{abcd}(\theta)
  =R_{ia}R_{jb}R_{kc}R_{ld}C^{(1)}_{ijkl}.
  \label{eq:cubic-tensor-rotation}
\end{equation}
The isotropic sagittal mode propagating along \(\mathbf n\) has displacement
\begin{equation}
  \widehat{\mathbf u}(z;\theta)
  =U(z)\mathbf n+W(z)\mathbf z,
  \label{eq:isotropic-sagittal-mode}
\end{equation}
and strain
\begin{align}
  \boldsymbol\epsilon
  ={}&E\,\mathbf n\mathbf n
   +Z\,\mathbf z\mathbf z
   +S(\mathbf n\mathbf z+\mathbf z\mathbf n),
  \notag\\
  &E=\mathrm{i}k_0U,
  \qquad
  Z=W',
  \qquad
  S=\frac{U'+\mathrm{i}k_0W}{2}.
  \label{eq:sagittal-strain-components}
\end{align}
Writing \(c=\cos\theta\) and \(s=\sin\theta\), its six independent
components in the crystal frame are
\begin{equation}
  \epsilon_{xx}=c^2E,
  \quad
  \epsilon_{yy}=s^2E,
  \quad
  \epsilon_{zz}=Z,
  \quad
  \epsilon_{xy}=csE,
  \quad
  \epsilon_{xz}=cS,
  \quad
  \epsilon_{yz}=sS.
  \label{eq:rotated-sagittal-strain}
\end{equation}

Before inserting the special family, the complete cubic strain contraction
is
\begin{align}
  \epsilon_{ij}^*C^{(1)}_{ijkl}\epsilon_{kl}
  ={}&C_{11}^{(1)}
      (|\epsilon_{xx}|^2+|\epsilon_{yy}|^2+|\epsilon_{zz}|^2)       \notag\\
    &+2C_{12}^{(1)}\operatorname{Re}
      (\epsilon_{xx}^*\epsilon_{yy}
       +\epsilon_{xx}^*\epsilon_{zz}
       +\epsilon_{yy}^*\epsilon_{zz})                              \notag\\
    &+4C_{44}^{(1)}
      (|\epsilon_{xy}|^2+|\epsilon_{xz}|^2+|\epsilon_{yz}|^2).
  \label{eq:complete-cubic-strain-contraction}
\end{align}
Every normal and shear component in
Eq.~\eqref{eq:rotated-sagittal-strain} is represented here.  The required
trigonometric products are
\begin{equation}
  c^4+s^4=\frac{3+\cos4\theta}{4},
  \qquad
  c^2s^2=\frac{1-\cos4\theta}{8}.
  \label{eq:fourfold-trigonometric-products}
\end{equation}
Substituting Eqs.~\eqref{eq:cubic-family-derivatives},
\eqref{eq:rotated-sagittal-strain}, and
\eqref{eq:fourfold-trigonometric-products} into
Eq.~\eqref{eq:complete-cubic-strain-contraction} gives, term by term,
\begin{align}
  \epsilon^*:C^{(1)}:\epsilon
  =\delta\left[
      \frac{5}{12}|E|^2
      +\frac{2}{3}|Z|^2
      -\frac{2}{3}\operatorname{Re}(E^*Z)
      +\frac{1}{4}|E|^2\cos4\theta
    \right].
  \label{eq:registered-cubic-strain-contraction}
\end{align}
The shear strain \(S\) has zero coefficient here precisely because
\(C_{44}^{(1)}=0\), not because it was excluded from the contraction.

For a mode of arbitrary mass norm
\begin{equation}
  N=\int_{-h}^{h}\rho(|U|^2+|W|^2)\,\mathrm dz,
  \label{eq:sagittal-mode-mass-norm}
\end{equation}
the first-order sensitivity is therefore
\begin{equation}
  V(\theta)=V_0+V_4\cos4\theta,
  \label{eq:first-order-cubic-angular-sensitivity}
\end{equation}
where
\begin{align}
  V_0
  &=\frac{\delta}{2\omega_0N}
    \int_{-h}^{h}\left[
      \frac{5}{12}|E|^2+\frac{2}{3}|Z|^2
      -\frac{2}{3}\operatorname{Re}(E^*Z)
    \right]\,\mathrm dz,                                           \label{eq:cubic-vzero}\\
  V_4
  &=Q_4\Delta_C^{(1)},
  \qquad
  Q_4:=\frac{1}{8\omega_0N}\int_{-h}^{h}|E|^2\,\mathrm dz
      =\frac{k_0^2}{8\omega_0N}\int_{-h}^{h}|U|^2\,\mathrm dz>0.
  \label{eq:cubic-vfour}
\end{align}
The last inequality holds for the selected finite-wavenumber Lamb mode,
whose in-plane component is nonzero.  Since
\(\Delta_C^{(1)}=\delta>0\) in the current family, it follows analytically
that \(V_4>0\).  The reference spectral evaluation gives
\(V_4=0.014063828174045103>0\) in the registered dimensionless units.
The physical angular frequency displacement is
\begin{equation}
  \varepsilon V_4
  =Q_4\Delta_C(\varepsilon),
  \label{eq:physical-cubic-angular-shift}
\end{equation}
so \(\varepsilon\) is applied exactly once.  Cubic symmetry by itself would
allow \(\cos(8\theta)\) and further multiples of four; the explicit
first-order tensor contraction proves that they are absent for this linear
perturbation family.

\subsection{Controlled normal form and radial implicit-function reduction}

We now state the local reduction with its derivative-level remainder.  The
regularity assumptions are stronger than the minimum needed for the radial
root because later stationary-phase arguments differentiate the phase and
amplitude further.

\begin{theorem}[Weak-anisotropy normal form]
\label{thm:anisotropic-normal-form}
Let the isotropic plate possess a positive-frequency ZGV circle at
\(k=k_0>0\), with radial curvature \(a\ne0\).  Assume that, in a closed
annulus about this circle and for \(|\varepsilon|\leq\varepsilon_0\), the
tracked generalized eigenvalue is simple, the mass matrix is positive
definite, and the mass-normalized eigengap obeys
\begin{equation}
  \inf_{\theta,q,\varepsilon}
  \min_{n\ne0}|\lambda_n-\lambda|=g>0.
  \label{eq:uniform-anisotropic-eigengap}
\end{equation}
Assume also that the Hermitian pencil is sufficiently smooth that the
selected real frequency branch is \(C^6\) in the local variables.  The
elastic pencil used here is analytic, so this hypothesis is automatic away
from an eigengap closure.  Then,
after reducing the annulus if necessary,
\begin{equation}
  \omega(q,\theta;\varepsilon)
  =\omega_0+\frac{a}{2}q^2
   +\varepsilon\bigl[V(\theta)+qB(\theta)\bigr]
   +R(q,\theta,\varepsilon),
  \label{eq:controlled-anisotropic-normal-form}
\end{equation}
where \(V=\partial_\varepsilon\omega(k_0,\theta;0)\),
\(B=\partial_k\partial_\varepsilon\omega(k_0,\theta;0)\), and
\begin{equation}
  R=q^3r_{30}+\varepsilon q^2r_{12}+\varepsilon^2r_{02},
  \label{eq:normal-form-remainder-factorization}
\end{equation}
with uniformly bounded periodic \(C^3\) coefficient functions.  In
particular, with \((x)_+=\max(x,0)\), every derivative used
in the two-dimensional Morse reduction satisfies
\begin{equation}
  \left|
    \partial_q^i\partial_\theta^jR(q,\theta,\varepsilon)
  \right|
  \leq C_{ij}\left(
    |q|^{3-i}
    +|\varepsilon|\,|q|^{(2-i)_+}
    +\varepsilon^2
  \right),
  \qquad 0\leq i+j\leq2.
  \label{eq:derivative-level-c2-remainder}
\end{equation}
The same bound with \(i=1\) holds after two angular derivatives, as follows
directly from Eq.~\eqref{eq:normal-form-remainder-factorization}.  There is a
unique periodic radial stationary graph
\(q=q_*(\theta,\varepsilon)\) near zero, and
\begin{equation}
  q_*(\theta,\varepsilon)
  =-\frac{\varepsilon B(\theta)}{a}
   +O_{C^2}(\varepsilon^2).
  \label{eq:radial-shift}
\end{equation}
If
\(
  W_2(\theta):=r_{02}(0,\theta,0)
  =\tfrac12\partial_\varepsilon^2
    \omega(k_0,\theta;0)
\), the frequency restricted to this graph is
\begin{equation}
  \omega_*(\theta;\varepsilon)
  :=\omega(q_*(\theta,\varepsilon),\theta;\varepsilon)
  =\omega_0+\varepsilon V(\theta)
   +\varepsilon^2\left[
      W_2(\theta)-\frac{B(\theta)^2}{2a}
    \right]
   +O_{C^2}(\varepsilon^3).
  \label{eq:reduced-angular-frequency-second-order}
\end{equation}
\end{theorem}

\paragraph{Proof.}
The uniform eigengap and simplicity make the spectral projector and the
selected eigenvalue smooth.  Taylor expansion first in \(q\) and then in
\(\varepsilon\), using
\(
  \omega_q(0,\theta;0)=0
\),
\(
  \omega_{qq}(0,\theta;0)=a
\), and isotropic independence of \(\theta\), gives
Eqs.~\eqref{eq:controlled-anisotropic-normal-form} and
\eqref{eq:normal-form-remainder-factorization}.  Applying the product rule
to the three factored terms gives
Eq.~\eqref{eq:derivative-level-c2-remainder} derivative by derivative.

Radial stationarity is the equation
\begin{equation}
  F(q,\theta,\varepsilon)
  :=\omega_q
  =aq+\varepsilon B(\theta)+R_q=0.
  \label{eq:radial-stationarity-equation}
\end{equation}
At \((q,\varepsilon)=(0,0)\), one has
\(F=0\) and \(F_q=a\ne0\).  The parameter-dependent
implicit-function theorem, applied uniformly on the compact angular circle,
gives the unique periodic graph.  It first gives \(q_*=O(\varepsilon)\).
The \(i=1\) bound in
Eq.~\eqref{eq:derivative-level-c2-remainder} then gives
\(
  R_q(q_*,\theta,\varepsilon)=O_{C^2}(\varepsilon^2)
\).  Solving Eq.~\eqref{eq:radial-stationarity-equation} for \(q_*\) proves
Eq.~\eqref{eq:radial-shift}.

At \(q_*=O(\varepsilon)\), the three parts of the remainder are
\begin{equation}
  q_*^3r_{30}=O_{C^2}(\varepsilon^3),
  \qquad
  \varepsilon q_*^2r_{12}=O_{C^2}(\varepsilon^3),
  \qquad
  \varepsilon^2r_{02}
  =\varepsilon^2W_2+O_{C^2}(\varepsilon^3).
  \label{eq:remainder-on-radial-graph}
\end{equation}
Finally,
\begin{equation}
  \frac a2q_*^2+\varepsilon q_*B
  =-\frac{\varepsilon^2B^2}{2a}
   +O_{C^2}(\varepsilon^3),
  \label{eq:radial-relaxation-energy}
\end{equation}
which proves Eq.~\eqref{eq:reduced-angular-frequency-second-order}.
\(\square\)

\subsection{Exact Cartesian Hessian and its Schur inertia}

The Morse classification must use the Cartesian Hessian, not the matrix of
uncorrected polar second partial derivatives.  For any scalar
\(f(q,\theta)\), with \(k=k_0+q>0\), direct differentiation of
\(
  \nabla_{\mathbf k}f
  =\mathbf n f_q+\mathbf t f_\theta/k
\)
gives the Cartesian Hessian in the orthonormal polar basis:
\begin{equation}
  \begin{pmatrix}\mathbf n&\mathbf t\end{pmatrix}^{\mathsf T}
  \nabla_{\mathbf k}^2f
  \begin{pmatrix}\mathbf n&\mathbf t\end{pmatrix}
  =
  \begin{pmatrix}
    f_{qq}
      &k^{-1}f_{q\theta}-k^{-2}f_\theta\\[2pt]
    k^{-1}f_{q\theta}-k^{-2}f_\theta
      &k^{-1}f_q+k^{-2}f_{\theta\theta}
  \end{pmatrix}.
  \label{eq:general-cartesian-polar-hessian}
\end{equation}
At a critical point \((q_j,\theta_j)\), both first partial derivatives
vanish, and the exact matrix is
\begin{equation}
  \mathbf H_j
  =
  \begin{pmatrix}
    \omega_{qq}
      &k_j^{-1}\omega_{q\theta}\\[2pt]
    k_j^{-1}\omega_{q\theta}
      &k_j^{-2}\omega_{\theta\theta}
  \end{pmatrix}_{(q_j,\theta_j)},
  \qquad k_j=k_0+q_j.
  \label{eq:critical-cartesian-polar-hessian}
\end{equation}

Along the radial graph, implicit differentiation gives
\begin{equation}
  q_*'=-\frac{\omega_{q\theta}}{\omega_{qq}},
  \qquad
  \omega_*'=\omega_\theta,
  \qquad
  \omega_*''
  =\omega_{\theta\theta}
   -\frac{\omega_{q\theta}^2}{\omega_{qq}}.
  \label{eq:reduced-angular-schur-derivative}
\end{equation}
If \(\omega_{qq}\ne0\), Sylvester's law of inertia applied to
Eq.~\eqref{eq:critical-cartesian-polar-hessian} gives the exact decomposition
\begin{equation}
  \operatorname{inertia}(\mathbf H_j)
  =\operatorname{inertia}(\omega_{qq})
   +\operatorname{inertia}\left(
      \frac{\omega_*''}{k_j^2}
    \right).
  \label{eq:hessian-schur-inertia}
\end{equation}
Thus the mixed derivative changes the tangential eigenvalue through the
Schur complement; it cannot be discarded before the inertia is taken.

At a critical point generated by the normal form,
Eqs.~\eqref{eq:derivative-level-c2-remainder} and
\eqref{eq:radial-shift} give
\begin{equation}
  \mathbf H_j
  =
  \begin{pmatrix}
    a+O(\varepsilon)
      &\dfrac{\varepsilon B'(\theta_j)}{k_0}+O(\varepsilon^2)\\[8pt]
    \dfrac{\varepsilon B'(\theta_j)}{k_0}+O(\varepsilon^2)
      &\dfrac{\varepsilon V''(\theta_j)}{k_0^2}
       +O(\varepsilon^2)
  \end{pmatrix}.
  \label{eq:leading-anisotropic-cartesian-hessian}
\end{equation}
Its Schur complement is
\begin{equation}
  \frac{1}{k_j^2}\left(
    \omega_{\theta\theta}
    -\frac{\omega_{q\theta}^2}{\omega_{qq}}
  \right)
  =\frac{\varepsilon V''(\theta_j)}{k_0^2}
   +O(\varepsilon^2).
  \label{eq:leading-tangential-schur-complement}
\end{equation}
For the reference branch \(a>0\), the radial inertia is positive and the
sign of \(\varepsilon V''\) determines whether the second direction is
positive or negative.

\subsection{General even harmonic and the cubic eight-point theorem}

Reciprocity gives
\(\omega(k,\theta;\varepsilon)
 =\omega(k,\theta+\pi;\varepsilon)\).
Therefore the order \(m\) of the first nonconstant angular harmonic must be
even.  A mirror axis fixes the phase origin of each allowed harmonic, but does
not by itself remove other harmonics.  We now impose the additional
pure-leading-harmonic hypothesis
\begin{equation}
  V(\theta)=V_0+V_m\cos[m(\theta-\theta_0)],
  \qquad V_m\ne0.
  \label{eq:pure-general-even-harmonic}
\end{equation}
Its derivative has the \(2m\) simple zeros
\begin{equation}
  \theta_j^{(0)}=\theta_0+\frac{j\pi}{m},
  \qquad j=0,\ldots,2m-1,
  \label{eq:pure-harmonic-angular-roots}
\end{equation}
and
\begin{equation}
  V''(\theta_j^{(0)})
  =-m^2V_m(-1)^j.
  \label{eq:pure-harmonic-angular-curvature}
\end{equation}
Since
\(
  \omega_*'=\varepsilon V'+O_{C^1}(\varepsilon^2)
\), the ordinary implicit-function theorem at each simple zero gives
\begin{equation}
  \theta_j=\theta_0+\frac{j\pi}{m}+O(\varepsilon),
  \qquad
  k_j=k_0-\frac{\varepsilon B(\theta_j)}{a}
       +O(\varepsilon^2).
  \label{eq:general-harmonic-critical-locations}
\end{equation}
For \(a>0\), Eq.~\eqref{eq:leading-tangential-schur-complement} shows that
the points alternate between \(m\) minima and \(m\) saddles.  Reversing the
sign of \(\varepsilon V_m\) leaves the leading angular locations fixed and
exchanges their roles.  Their leading frequencies are
\begin{equation}
  \omega_j
  =\omega_0+\varepsilon[V_0+V_m(-1)^j]
   +O(\varepsilon^2),
  \qquad
  \Delta\omega_{\mathrm{ms}}
  =2|\varepsilon V_m|+O(\varepsilon^2).
  \label{eq:general-harmonic-frequency-splitting}
\end{equation}
The gradient index is \(+1\) at each minimum and \(-1\) at each saddle, so
the algebraic index sum is
\begin{equation}
  m(+1)+m(-1)=0.
  \label{eq:general-harmonic-index-sum}
\end{equation}

A fixed higher harmonic must be distinguished from the
\(O(\varepsilon)\) correction to the roots above.  Suppose instead that
\begin{equation}
  V(\theta)=V_0+V_m\cos[m(\theta-\theta_0)]
             +\mathcal H(\theta).
  \label{eq:leading-harmonic-plus-remainder}
\end{equation}
There is a constant \(c_m>0\) such that the dominance condition
\begin{equation}
  \lVert\mathcal H\rVert_{C^2}
  <c_m|V_m|
  \label{eq:higher-harmonic-dominance}
\end{equation}
preserves the \(2m\) simple zeros and their alternating curvatures.  Their
locations obey the separate estimate
\begin{equation}
  \theta_j
  =\theta_0+\frac{j\pi}{m}
   +O\left(
      \frac{\lVert\mathcal H\rVert_{C^1}}{|V_m|}
    \right)
   +O(\varepsilon).
  \label{eq:higher-harmonic-root-shift}
\end{equation}
Without a bound such as Eq.~\eqref{eq:higher-harmonic-dominance}, point-group
symmetry alone does not fix the number of angular critical points.

\begin{theorem}[Cubic Morse splitting]
\label{thm:cubic-morse-splitting}
Assume the hypotheses of Theorem~\ref{thm:anisotropic-normal-form}, take the
reference radial curvature \(a>0\), and use the [001] cubic family
\eqref{eq:registered-cubic-family} with \(\delta>0\).  Then, for every
sufficiently small nonzero \(\varepsilon\), a thin annular neighbourhood of
the isotropic ZGV circle contains exactly eight critical points of the
tracked branch.  They satisfy
\begin{equation}
  \theta_j=\frac{j\pi}{4}+O(\varepsilon),
  \qquad
  k_j=k_0-\frac{\varepsilon B(\theta_j)}{a}
       +O(\varepsilon^2),
  \qquad j=0,\ldots,7.
  \label{eq:cubic-eight-critical-points}
\end{equation}
Four are Morse minima and four are Morse saddles.  For
\(\varepsilon>0\), the odd \(j\) are minima and the even \(j\) are saddles;
for \(\varepsilon<0\), these roles reverse.  The two critical-frequency
families are separated by
\begin{equation}
  \Delta\omega_{\mathrm{ms}}
  =2|\varepsilon V_4|+O(\varepsilon^2),
  \qquad V_4=Q_4\delta>0,
  \label{eq:cubic-minimum-saddle-separation}
\end{equation}
and their total gradient index is zero.
\end{theorem}

\paragraph{Proof.}
Equations~\eqref{eq:registered-cubic-strain-contraction}--
\eqref{eq:cubic-vfour} prove, rather than infer from symmetry, that the
first-order angular sensitivity is the pure harmonic
\(V_0+V_4\cos4\theta\) and that \(V_4>0\).  Set \(m=4\) and
\(\theta_0=0\) in
Eqs.~\eqref{eq:pure-harmonic-angular-roots}--
\eqref{eq:general-harmonic-critical-locations}.  This gives eight points.
At their leading locations,
\begin{equation}
  \varepsilon V''(\theta_j)
  =-16\varepsilon V_4(-1)^j+O(\varepsilon^2).
  \label{eq:cubic-tangential-curvature-sign}
\end{equation}
The radial Schur pivot is positive, while the second pivot has the sign in
Eq.~\eqref{eq:cubic-tangential-curvature-sign}.  The stated minimum--saddle
classification and its reversal follow.  Equation
\eqref{eq:general-harmonic-frequency-splitting} supplies the separation, and
Eq.~\eqref{eq:general-harmonic-index-sum} supplies the index sum.
To prove that the count is exact, choose disjoint angular neighbourhoods of
the eight simple zeros of \(V'\).  On their compact complement,
\(\lvert V'\rvert\) has a positive lower bound.  After division by the
nonzero \(\varepsilon\), the \(O_{C^1}(\varepsilon)\) correction in
\(\omega_*'/\varepsilon\) cannot create a zero there when
\(\lvert\varepsilon\rvert\) is sufficiently small.  The radial graph is
unique at every angle by Theorem~\ref{thm:anisotropic-normal-form}; hence no
additional critical point exists in the reduced annulus.
\(\square\)

\subsection{Exceptional cases and theorem boundaries}

The condition \(V_m\ne0\) is substantive.  If a nominal coefficient such
as \(V_4\) vanishes, the following cases must be distinguished.
\begin{enumerate}
  \item If a higher harmonic \(V_\ell\cos(\ell\theta)\), with even
        \(\ell>m\), is the first nonzero nonconstant term at order
        \(\varepsilon\), the preceding calculation applies with
        \(m\) replaced by \(\ell\).
  \item If the whole first-order sensitivity is angularly constant, the
        second order reduced coefficient is
        \begin{equation}
          V_2^{\mathrm{eff}}(\theta)
          :=W_2(\theta)-\frac{B(\theta)^2}{2a}.
          \label{eq:effective-second-order-angular-potential}
        \end{equation}
        A nonzero nondegenerate harmonic of
        \(V_2^{\mathrm{eff}}\) splits the ring with tangential curvature of
        order \(\varepsilon^2\), not order \(\varepsilon\).
  \item If two or more harmonics compete at the same first nonzero order,
        the critical-point count is the zero count of their full derivative;
        it is not determined by one nominal order.  If at a root both
        \(V'=0\) and \(V''=0\), the point is degenerate and the Morse
        conclusion does not apply until the next local term is resolved.
\end{enumerate}

Two independent spectral boundaries are equally important.  If \(a=0\),
then \(F_q(0,\theta,0)=0\), the radial implicit-function argument fails,
and the balance between anisotropy and the first nonzero higher radial term
has a different scaling.  If the eigengap \(g\) closes, the reduced
resolvent becomes unbounded and a scalar tracked branch is not smooth; a
matrix-valued degenerate perturbation problem is then required.  The cases
\(k_0=0\) and \(\omega_0=0\) are also outside the present polar-ring and
positive-square-root formulation, respectively.

\subsection{Conversion to the dimensionless computational convention}

The code uses
\begin{equation}
  \kappa=kh,
  \qquad
  \Omega=\frac{\omega h}{c_T},
  \qquad
  c_T=\sqrt{\frac{\mu}{\rho}},
  \qquad
  q_\kappa=\kappa-\kappa_0=hq.
  \label{eq:anisotropic-nondimensional-variables}
\end{equation}
Writing subscripts \(\omega\) and \(\Omega\) for coefficients in the
dimensional and dimensionless frequency normal forms, the chain rule gives
\begin{equation}
  V_\Omega=\frac{h}{c_T}V_\omega,
  \qquad
  B_\Omega=\frac{1}{c_T}B_\omega,
  \qquad
  a_\Omega=\frac{1}{hc_T}a_\omega.
  \label{eq:anisotropic-coefficient-conversion}
\end{equation}
Hence
\begin{equation}
  q_{\kappa,*}
  =-\frac{\varepsilon B_\Omega}{a_\Omega}
   +O(\varepsilon^2)
  =h\left[-\frac{\varepsilon B_\omega}{a_\omega}
   +O(\varepsilon^2)\right],
  \label{eq:dimensionless-radial-shift}
\end{equation}
which is exactly the dimensional location multiplied by \(h\).  If the
dimensional invariant \(\Delta_C^{(1)}\) is retained, then
\(Q_{4,\Omega}=(h/c_T)Q_{4,\omega}\).  If stiffness is additionally divided
by \(\mu\), the corresponding coefficient absorbs the factor \(\mu\).
For the registered reference values
\(h=\rho=\mu=1\), one has \(c_T=1\), so the numerical and dimensional
coefficient values coincide in the chosen units.
